\section{Proofs for the stochastic theory}\label{app:stoch-proofs}

We prove the statements of \Cref{sec:stochastic} deferred there. Throughout
we use the notation of that section: $g(a)=-\E_{\rho_a}[\nabla V]$, $A(a)=\E_{\rho_a}[\nabla^2V]$,
$R(a)=C^{-1}-A(a)$, the intrinsic fluctuation $\Psi$ of \Cref{def:Psi}, and the minibatch estimators
\eqref{eq:STL-score}.

\subsection{The rescued initializer}

\begin{proof}[Proof of \Cref{lem:rescue-energy}]
Write $d:=m_{\mathrm{in}}-m_\star$ and $\bar H:=\int_0^1\nabla^2V(m_\star+td)\dd t$, so that
$\nabla V(m_{\mathrm{in}})=\nabla V(m_\star)+\bar Hd$ by the integral Taylor form. Note that $m_\star$
is the optimal variational mean, not the mode of the target: the stationarity condition is
$\E_{\rho_{a_\star}}[\nabla V]=0$, which does not force $\nabla V(m_\star)=0$.
We first bound $\nabla V(m_\star)$. With $X=m_\star+C_\star^{1/2}Z$ and Taylor expansion,
$\nabla V(X)=\nabla V(m_\star)+\nabla^2V(m_\star)C_\star^{1/2}Z+r(X)$ with
$\norm{r(X)}\leq\tfrac{\LH}2\norm{C_\star^{1/2}Z}^2$. Taking $\E_{\rho_{a_\star}}$, the linear term
vanishes and $\E_{\rho_{a_\star}}[\nabla V]=0$, so
\begin{equation}
\label{eq:gradV-mstar}
    \norm{\nabla V(m_\star)}\leq\tfrac{\LH}2\E\norm{C_\star^{1/2}Z}^2=\tfrac{\LH}2\Tr C_\star\leq\frac{\LH N_\theta}{2\alpha}.
\end{equation}
For the mean, since $C_0^{-1}=H_{\mathrm{in}}$,
\[
    m_0-m_\star=d-H_{\mathrm{in}}^{-1}\bigl(\nabla V(m_\star)+\bar Hd\bigr)
    =H_{\mathrm{in}}^{-1}(H_{\mathrm{in}}-\bar H)d-H_{\mathrm{in}}^{-1}\nabla V(m_\star).
\]
Hessian-Lipschitzness gives $\norm{H_{\mathrm{in}}-\bar H}\leq\int_0^1\LH(1-t)r_{\mathrm{in}}\dd t=\tfrac{\LH}2r_{\mathrm{in}}$,
and $\norm{H_{\mathrm{in}}^{-1}}\leq1/\alpha$, so with \eqref{eq:gradV-mstar},
$\norm{m_0-m_\star}\leq\tfrac1\alpha\cdot\tfrac{\LH}2r_{\mathrm{in}}^2+\tfrac1\alpha\cdot\tfrac{\LH N_\theta}{2\alpha}
=\tfrac{\LH}{2\alpha}(r_{\mathrm{in}}^2+\tfrac{N_\theta}\alpha)$.
For the covariance, with $C_\star^{-1}=A(a_\star)=\E_{\rho_{a_\star}}[\nabla^2V]$,
\[
    \norm{H_{\mathrm{in}}-C_\star^{-1}}\leq\E_{\rho_{a_\star}}\norm{\nabla^2V(m_{\mathrm{in}})-\nabla^2V(X)}
    \leq\LH\,\E_{\rho_{a_\star}}\norm{m_{\mathrm{in}}-X}\leq\LH\bigl(r_{\mathrm{in}}+\sqrt{N_\theta/\alpha}\bigr),
\]
using $\E_{\rho_{a_\star}}\norm{m_{\mathrm{in}}-X}\leq r_{\mathrm{in}}+\sqrt{\Tr C_\star}$ and
$\Tr C_\star\leq N_\theta/\alpha$. From
$C_0-C_\star=H_{\mathrm{in}}^{-1}(C_\star^{-1}-H_{\mathrm{in}})C_\star$ and
$\norm X_{\Frob}\leq\sqrt{N_\theta}\norm X_\op$,
$\norm{C_0-C_\star}_{\Frob}\leq\tfrac{\sqrt{N_\theta}\LH}{\alpha^2}(r_{\mathrm{in}}+\sqrt{N_\theta/\alpha})$.
It remains to bound the energy gap. Since $a_\star$ is stationary, the gap is second order in the
parameter error.
Parameterize $f(m,B):=\E_Z[V(m+BZ)]-\log\det B$, so $f(m,B)=\calE(m,BB^\top)+\text{const}$ and
$\Delta(a_0)=f(m_0,B_0)-f(m_\star,B_\star)$ with $B_0:=C_0^{1/2}$, $B_\star:=C_\star^{1/2}$. Along the
affine path $m_s=m_\star+s\,\delta m$, $B_s=B_\star+s\,\delta B$ with $\delta m=m_0-m_\star$ and
$\delta B=B_0-B_\star$ symmetric, the first derivative vanishes at $s=0$: $\partial_mf=\E[\nabla V]=0$
at $a_\star$, and by Stein's identity
$\tfrac{\dd}{\dd s}\E[V(m_\star+B_sZ)]\big|_0=\Tr(\delta B\,\E[\nabla^2V(X)]B_\star)=\Tr(\delta B\,B_\star^{-1})
=\tfrac{\dd}{\dd s}\log\det B_s\big|_0$. For the second derivatives,
\[
    \frac{\dd^2}{\dd s^2}\E[V(m_s+B_sZ)]
    =\E\bigl[(\delta m+\delta BZ)^\top\nabla^2V(X_s)(\delta m+\delta BZ)\bigr]
    \leq\beta\bigl(\norm{\delta m}^2+\norm{\delta B}_{\Frob}^2\bigr),
\]
the cross term vanishing by $\E Z=0$, and with $M_s:=B_s^{-1}\delta B$,
$\tfrac{\dd^2}{\dd s^2}[-\log\det B_s]=\Tr(M_s^2)\leq\norm{M_s}_{\Frob}^2\leq\beta\norm{\delta B}_{\Frob}^2$,
using $\Tr(M^2)\leq\norm M_{\Frob}^2$ and $\norm{B_s^{-1}}_\op\leq\sqrt\beta$, both $C_0$ and $C_\star$
lying on the band $\succeq\tfrac1\beta\Id$. Taylor with vanishing first derivative gives
$\Delta(a_0)\leq\tfrac\beta2\norm{m_0-m_\star}^2+\beta\norm{C_0^{1/2}-C_\star^{1/2}}_{\Frob}^2$, and the
square root is Frobenius-Lipschitz on the band, eigenvalue divided differences of $\sqrt\cdot$ on
$[\tfrac1\beta,\infty)$ being at most $\tfrac{\sqrt\beta}2$, so
$\norm{C_0^{1/2}-C_\star^{1/2}}_{\Frob}\leq\tfrac{\sqrt\beta}2\norm{C_0-C_\star}_{\Frob}$. Substituting
the two parameter bounds gives the stated estimate; every term carries $\LH^2$, so $\Delta(a_0)=0$ when
$\LH=0$.
\end{proof}

\subsection{The global one-step estimates}

\begin{proof}[Proof of \Cref{lem:onestep-stoch-Riem}]
By \Cref{lem:stoch-band}(1) the band $C_n\preceq\tfrac1\alpha\Id$ holds pathwise, so the scope
condition $\norm{C_n^{-1/2}UC_n^{-1/2}}_\op\leq1$ of the retraction-smoothness estimate
\eqref{eq:retr-smooth-restated} holds pathwise for the tangent
$\zeta_n=-\Delta t\,\widehat\grad\calE(a_n)$ because $\Delta t\kappa\leq1$, and
\eqref{eq:retr-smooth-restated} with $\Lambda=\tfrac1\alpha$, hence $L=2\beta\Lambda=2\kappa$, gives
pathwise
\[
    \calE(a_{n+1})\leq\calE(a_n)+\inner{\grad\calE(a_n)}{\zeta_n}_{a_n}+\kappa\Delta t^2\norm{\widehat\grad\calE(a_n)}_{a_n}^2 .
\]
Taking $\E[\cdot\mid\Fil_n]$, conditional unbiasedness \eqref{eq:STL-unbiased} gives
$\E[\inner{\grad\calE(a_n)}{\widehat\grad\calE(a_n)}_{a_n}\mid\Fil_n]=\Gradnorm(a_n)^2$, and the
intrinsic variance bound \eqref{eq:intrinsic-STL} gives
$\E[\norm{\widehat\grad\calE(a_n)}_{a_n}^2\mid\Fil_n]\leq(1+\tfrac2B)\Gradnorm(a_n)^2+\tfrac3{2B}\Psi(a_n)$.
Collecting terms gives \eqref{eq:Riem-onestep-Psi}.
\end{proof}

\begin{proof}[Proof of \Cref{lem:onestep-stoch-KL}]
Suppress the index $n$ and condition on $\Fil_n$. Write $A:=A(a)$, $R:=C^{-1}-A$, $S:=C^{1/2}RC^{1/2}$,
$P_0:=C^{1/2}AC^{1/2}=\Id-S$, and $D:=P-P_0$, so that $\overline S=S-D$ with $\E[D\mid\Fil]=0$ and
$\E[\norm D_{\Frob}^2\mid\Fil]=\Psi(a)/B$ by \Cref{def:Psi} and minibatch independence. Set
\[
    M:=g(a)^\top C\,g(a),\qquad V:=\norm S_{\Frob}^2,\qquad\sigma^2:=\frac{\Psi(a)}B,
    \qquad\text{so that}\qquad\Gradnorm(a)^2=M+\tfrac12V .
\]
On the band $C\preceq\tfrac1\alpha\Id$ the retraction-smoothness estimate
\eqref{eq:retr-smooth-restated} applies with $\Lambda=\tfrac1\alpha$ and $L=2\kappa$, and by
\Cref{lem:KL-retraction} the tangent $\zeta=(\Delta t\,C\overline b,\,C^{1/2}W(P)C^{1/2})$ meets its
scope condition, so pathwise
$\calE(a^+)\leq\calE(a)+\inner{\grad\calE(a)}\zeta_a+\kappa\norm\zeta_a^2$. The first-order term
evaluates to $\inner{\grad\calE(a)}\zeta_a=-\Delta t\,g^\top C\overline b-\tfrac12\Tr(S\,W(P))$: the
mean block uses $\grad_m\calE=-Cg$ in the metric $v^\top C^{-1}v$, and the covariance block uses
$\grad_C\calE=-CRC$ in the metric $\tfrac12\Tr(C^{-1}\cdot C^{-1}\cdot)$, giving
$\tfrac12\Tr(C^{-1}(-CRC)C^{-1}C^{1/2}W(P)C^{1/2})=-\tfrac12\Tr(S\,W(P))$. Likewise
$\norm\zeta_a^2=\Delta t^2\,\overline b^\top C\overline b+\tfrac12\norm{W(P)}_{\Frob}^2$. Hence,
pathwise,
\[
    \calE(a^+)-\calE(a)
    \leq-\Delta t\,g^\top C\overline b-\tfrac12\Tr(S\,W(P))
    +\kappa\,\Delta t^2\,\overline b^\top C\overline b+\tfrac\kappa2\norm{W(P)}_{\Frob}^2 .
\]
Write $\overline b=g+\varepsilon_b$ with $\E[\varepsilon_b\mid\Fil]=0$. By
\eqref{eq:intrinsic-STL} and independence,
$\E[\varepsilon_b^\top C\varepsilon_b\mid\Fil]\leq(V+\Psi(a))/B=\tfrac VB+\sigma^2$, so
$\E[\overline b^\top C\overline b\mid\Fil]\leq M+\tfrac VB+\sigma^2$ and
\begin{equation}
\label{eq:KL-mean}
    \E\bigl[-\Delta t\,g^\top C\overline b+\kappa\,\Delta t^2\,\overline b^\top C\overline b\mid\Fil\bigr]
    \leq-\Delta t(1-\kappa\Delta t)M+\frac{\kappa\,\Delta t^2}BV+\kappa\,\Delta t^2\sigma^2 .
\end{equation}
For the covariance drift, with $W(P)=\Delta t\,\overline S+\mathcal R(P)$ we have
$-\tfrac12\Tr(S\,W(P))=-\tfrac{\Delta t}2\Tr(S\overline S)-\tfrac12\Tr(S\,\mathcal R(P))$, and
$\E[\overline S\mid\Fil]=S$ gives $\E[\Tr(S\overline S)\mid\Fil]=V$. For the remainder,
Cauchy--Schwarz, \eqref{eq:Wlog-bounds}, $\E[\norm{\overline S}_{\Frob}^2\mid\Fil]=V+\sigma^2$,
Jensen, and $x\sqrt{x^2+y^2}\leq x^2+\tfrac12y^2$ with $x=\sqrt V$, $y=\sigma$ give
$|\E[\Tr(S\,\mathcal R(P))\mid\Fil]|\leq\kappa\Delta t^2(V+\tfrac12\sigma^2)$, so
\begin{equation}
\label{eq:KL-covdrift}
    \E\bigl[-\tfrac12\Tr(S\,W(P))\mid\Fil\bigr]
    \leq-\frac{\Delta t}2V+\frac{\kappa\,\Delta t^2}2V+\frac{\kappa\,\Delta t^2}4\sigma^2 .
\end{equation}
Likewise, from $\norm{W(P)}_{\Frob}\leq\Delta t\norm{\overline S}_{\Frob}$ we get
$\E[\norm{W(P)}_{\Frob}^2\mid\Fil]\leq\Delta t^2(V+\sigma^2)$, so
\begin{equation}
\label{eq:KL-covquad}
    \frac\kappa2\E[\norm{W(P)}_{\Frob}^2\mid\Fil]\leq\frac{\kappa\,\Delta t^2}2V+\frac{\kappa\,\Delta t^2}2\sigma^2 .
\end{equation}
Adding \eqref{eq:KL-mean}, \eqref{eq:KL-covdrift} and \eqref{eq:KL-covquad} to
$\Delta_n$,
\[
    \E[\Delta_{n+1}\mid\Fil]
    \leq\Delta_n-\Delta t(1-\kappa\Delta t)M
    -\Delta t\Bigl[\tfrac12-\kappa\Delta t\bigl(1+\tfrac1B\bigr)\Bigr]V
    +\tfrac74\,\kappa\,\Delta t^2\sigma^2 .
\]
For $\Delta t\leq1/(8\kappa)$ one has $1-\kappa\Delta t\geq\tfrac78\geq\tfrac12$ and, since $B\geq1$,
$\tfrac12-\kappa\Delta t(1+\tfrac1B)\geq\tfrac12-\tfrac14=\tfrac14$. Thus
$\E[\Delta_{n+1}\mid\Fil]\leq\Delta_n-\tfrac{\Delta t}2M-\tfrac{\Delta t}4V+\tfrac{7\kappa\Delta t^2}{4B}\Psi(a)$,
which is \eqref{eq:KL-onestep-Psi} since $\Gradnorm(a)^2=M+\tfrac12V$.
\end{proof}

\subsection{The global theorems}

\begin{proof}[Proof of \Cref{thm:stoch-glob-Riem}]
\Cref{lem:stoch-band}(1) gives the band $\tfrac1{2\beta}\Id\preceq C_n\preceq\tfrac1\alpha\Id$
pathwise, so the one-step estimate \eqref{eq:Riem-onestep-Psi} and the gradient domination
$\Gradnorm(a_n)^2\geq\tfrac1{2\kappa}\Delta_n$ of \eqref{eq:PL-bands} are available at every step.
For (1), at $\Delta t=\tfrac1{2\kappa(1+2/B)}$ the bracket in \eqref{eq:Riem-onestep-Psi} equals
$\tfrac12$, so with $\Psi\leq\Psi_H=N_\theta^2\barLH^2$ from \Cref{lem:Psi-hesslip},
$\E[\Delta_{n+1}\mid\Fil_n]\leq(1-\tfrac{\Delta t}{4\kappa})\Delta_n+\tfrac{3\kappa\Delta t^2}{2B}\Psi_H$.
Unrolling, the floor is
$\tfrac{3\kappa\Delta t^2\Psi_H/(2B)}{\Delta t/(4\kappa)}=\tfrac{6\kappa^2\Delta t}B\Psi_H=\tfrac{3\kappa}{B+2}\Psi_H$
and the contraction is $1-\tfrac{\Delta t}{4\kappa}=1-\tfrac1{8\kappa^2(1+2/B)}$.
For (2), on the band $C_n\preceq\tfrac1\alpha\Id$ we have $\beta\Lambda=\kappa$, so
\Cref{lem:Psi-selfbound} gives $\tfrac32\Psi(a_n)\leq3\kappa N_\theta+3\kappa\Gradnorm(a_n)^2$.
Substituting into \eqref{eq:Riem-onestep-Psi} and collecting the $\Gradnorm^2$ terms,
\[
    \E[\Delta_{n+1}\mid\Fil_n]
    \leq\Delta_n-\Delta t\bigl(1-\kappa\Delta t\rho_B\bigr)\Gradnorm(a_n)^2+\frac{3\kappa^2\Delta t^2N_\theta}B,
\]
since $\kappa\Delta t^2(1+\tfrac2B)+\tfrac{3\kappa\Delta t^2}{2B}\cdot2\kappa=\kappa\Delta t^2\rho_B$
with $\rho_B=1+\tfrac{2+3\kappa}B$. At $\Delta t=\tfrac1{2\kappa\rho_B}$ the bracket is $\tfrac12$, so
$\E[\Delta_{n+1}\mid\Fil_n]\leq(1-\tfrac{\Delta t}{4\kappa})\Delta_n+\tfrac{3\kappa^2\Delta t^2N_\theta}B$.
Unrolling, the floor is
$\tfrac{3\kappa^2\Delta t^2N_\theta/B}{\Delta t/(4\kappa)}=\tfrac{12\kappa^3\Delta tN_\theta}B=\tfrac{6\kappa^2N_\theta}{B\rho_B}=\tfrac{6\kappa^2N_\theta}{B+2+3\kappa}$
and the contraction is $1-\tfrac1{8\kappa^2\rho_B}$. Since $B+2+3\kappa\geq3\kappa$, the floor is at
most $2\kappa N_\theta$.
\end{proof}

\begin{proof}[Proof of \Cref{thm:stoch-glob-KL}]
\Cref{lem:stoch-band}(2) gives the band $\tfrac1\beta\Id\preceq C_n\preceq\tfrac1\alpha\Id$ pathwise
for every $\Delta t>0$, and \eqref{eq:PL-bands} gives $\Gradnorm(a_n)^2\geq\tfrac1\kappa\Delta_n$.
For (1), by \Cref{lem:Psi-hesslip}, $\Psi\leq\Psi_H=N_\theta^2\barLH^2$; substituting into
\eqref{eq:KL-onestep-Psi} gives
$\E[\Delta_{n+1}\mid\Fil_n]\leq(1-\tfrac{\Delta t}{2\kappa})\Delta_n+\tfrac{7\kappa\Delta t^2}{4B}\Psi_H$.
At $\Delta t=\tfrac1{8\kappa}$ the contraction is $1-\tfrac1{16\kappa^2}$ and the floor is
$\tfrac{7\kappa\Delta t^2\Psi_H/(4B)}{\Delta t/(2\kappa)}=\tfrac{7\kappa}{16B}\Psi_H$.
For (2), we have $\beta\Lambda=\kappa$ on the band, so \Cref{lem:Psi-selfbound} gives
$\Psi(a_n)\leq2\kappa N_\theta+2\kappa\Gradnorm(a_n)^2$, and \eqref{eq:KL-onestep-Psi} becomes
\[
    \E[\Delta_{n+1}\mid\Fil_n]
    \leq\Delta_n-\frac{\Delta t}2\Bigl(1-\frac{7\kappa^2\Delta t}B\Bigr)\Gradnorm(a_n)^2+\frac{7\kappa^2\Delta t^2}{2B}N_\theta .
\]
The choice $\Delta t\leq\tfrac B{14\kappa^2}$ makes $7\kappa^2\Delta t/B\leq\tfrac12$, so the bracket is
at least $\tfrac12$, and gradient domination gives
$\E[\Delta_{n+1}\mid\Fil_n]\leq(1-\tfrac{\Delta t}{4\kappa})\Delta_n+\tfrac{7\kappa^2\Delta t^2}{2B}N_\theta$.
Unrolling, the floor is
$\tfrac{7\kappa^2\Delta t^2N_\theta/(2B)}{\Delta t/(4\kappa)}=\tfrac{14\kappa^3\Delta tN_\theta}B$,
which equals $\kappa N_\theta$ when $\Delta t=\tfrac B{14\kappa^2}$, and the contraction is then
$1-\tfrac B{56\kappa^3}$.
\end{proof}

\subsection{Local variance and the matrix-exponential perturbation}

\begin{proof}[Proof of \Cref{lem:loc-variance}]
Write $H(x):=\nabla^2\widehat V(x)$, so $\varepsilon_H=H(X)-A(a)$ and $A(a)=\E_{\rho_a}[H(X)]$.
Gaussian Poincar\'e applied componentwise to $X\mapsto H_{ij}(X)$ with $X\sim\N(m,C)$
gives $\Var(H_{ij}(X))\leq\norm C_\op\E\norm{\nabla H_{ij}(X)}^2$, so
$\E\norm{\varepsilon_H}_{\Frob}^2\leq\norm C_\op\E\sum_{i,j,k}(\partial_kH_{ij}(X))^2$. For each $k$,
$\sum_{i,j}(\partial_kH_{ij}(x))^2=\norm{DH(x)[e_k]}_{\Frob}^2\leq N_\theta\norm{DH(x)[e_k]}_\op^2\leq N_\theta\LHs^2$
by \eqref{eq:LHs}; summing over the $N_\theta$ coordinates $k$ and using $\norm C_\op\leq\tfrac32$,
\begin{equation}
\label{eq:covnoise}
    \E\norm{\varepsilon_H}_{\Frob}^2\leq\tfrac32N_\theta^2\LHs^2 .
\end{equation}
Let $F(x):=-\nabla\widehat V(x)+C^{-1}(x-m)$, so $\varepsilon_b=F(X)-\E F(X)$ and
$DF(x)=C^{-1}-H(x)$ is symmetric. Gaussian Poincar\'e componentwise and summed gives
$\E\norm{\varepsilon_b}^2\leq\norm C_\op\E\norm{C^{-1}-H(X)}_{\Frob}^2$, and since $\E H(X)=A(a)$ the
orthogonal decomposition $C^{-1}-H(X)=(C^{-1}-A(a))-\varepsilon_H$ has vanishing cross expectation, so
\begin{equation}
\label{eq:meannoise}
    \E\norm{\varepsilon_b}^2\leq\tfrac32\Bigl(\norm{C^{-1}-A(a)}_{\Frob}^2+\E\norm{\varepsilon_H}_{\Frob}^2\Bigr).
\end{equation}
Adding $\tfrac12\E\norm{\varepsilon_H}_{\Frob}^2$ to \eqref{eq:meannoise},
\[
    \E\Bigl[\norm{\varepsilon_b}^2+\tfrac12\norm{\varepsilon_H}_{\Frob}^2\Bigr]
    \leq2\,\E\norm{\varepsilon_H}_{\Frob}^2+\tfrac32\norm{C^{-1}-A(a)}_{\Frob}^2 ,
\]
and by \eqref{eq:covnoise} the first term is at most $3N_\theta^2\LHs^2=V_0$. For the second,
$\norm{C^{-1}-A(a)}_{\Frob}^2\leq2\norm{C^{-1}-\Id}_{\Frob}^2+2\norm{A(a)-\Id}_{\Frob}^2$; since
$C^{-1}-\Id=-C^{-1}(C-\Id)$ and $\norm{C^{-1}}_\op\leq2$, the first piece is at most
$8\norm{C-\Id}_{\Frob}^2$. For the second piece, $A(a)-\Id=\E_Z[H(m+C^{1/2}Z)-H(Z)]$ because
$A(a_\star)=\E_ZH(Z)=\Id$, so by Jensen and \eqref{eq:LHs},
\[
    \norm{A(a)-\Id}_{\Frob}^2\leq N_\theta\LHs^2\,\E\norm{m+(C^{1/2}-\Id)Z}^2
    =N_\theta\LHs^2\bigl(\norm m^2+\norm{C^{1/2}-\Id}_{\Frob}^2\bigr),
\]
the cross term vanishing by $\E Z=0$, and $\norm{C^{1/2}-\Id}_{\Frob}^2\leq\norm{C-\Id}_{\Frob}^2$
eigenvalue-wise. Hence
$\tfrac32\norm{C^{-1}-A(a)}_{\Frob}^2\leq12\norm{C-\Id}_{\Frob}^2+3N_\theta\LHs^2(\norm m^2+\norm{C-\Id}_{\Frob}^2)$.
Converting through $\norm{a-a_\star}_\star^2=\norm m^2+\tfrac12\norm{C-\Id}_{\Frob}^2$ gives
$\norm{C-\Id}_{\Frob}^2\leq2\norm{a-a_\star}_\star^2$ and
$\norm m^2+\norm{C-\Id}_{\Frob}^2\leq2\norm{a-a_\star}_\star^2$, so the state-dependent part is at most
$(24+6N_\theta\LHs^2)\norm{a-a_\star}_\star^2=V_1\norm{a-a_\star}_\star^2$. The minibatch bound is
$\E\norm{\overline\varepsilon}^2=B^{-1}\E\norm\varepsilon^2$ for centered independent averages. At
$a_\star$ we have $C=\Id$ and $C^{-1}-A(a_\star)=0$, so \eqref{eq:covnoise} tightens to
$\E\norm{\varepsilon_H}_{\Frob}^2\leq N_\theta^2\LHs^2$ and \eqref{eq:meannoise} to
$\E\norm{\varepsilon_b}^2\leq N_\theta^2\LHs^2$, giving $\sigma_\star^2\leq\tfrac32N_\theta^2\LHs^2$.
\end{proof}

The only nonlinear ingredient of the local perturbation analysis is the following pair of elementary
bounds on the matrix exponential.

\begin{lemma}[Exponential perturbation bounds]
\label{lem:exp-perturb}
Let $S,\widehat S\in\Sym(N_\theta)$ with $S\preceq\Id$ and $\widehat S\preceq\Id$, put
$\Xi:=\widehat S-S$, and let $0<\Delta t\leq1$. Then
\begin{align*}
\textup{(i)}\quad & \norm{e^{\Delta t\widehat S}-e^{\Delta tS}}_{\Frob}\leq e\,\Delta t\,\norm{\Xi}_{\Frob},\\
\textup{(ii)}\quad & \norm{e^{\Delta t\widehat S}-e^{\Delta tS}-D\exp(\Delta tS)[\Delta t\,\Xi]}_{\Frob}
\leq\tfrac e2\,\Delta t^2\,\norm\Xi_\op\norm\Xi_{\Frob}.
\end{align*}
\end{lemma}

\begin{proof}
For symmetric $M=Q\diag(\lambda_i)Q^\top$ the Fr\'echet derivative acts in the eigenbasis by the
Loewner divided-difference matrix,
$(Q^\top D\exp(M)[E]Q)_{ij}=\Phi_{ij}(Q^\top EQ)_{ij}$ with
$\Phi_{ij}=(e^{\lambda_i}-e^{\lambda_j})/(\lambda_i-\lambda_j)$, equal to $e^{\lambda_i}$ when
$\lambda_i=\lambda_j$; since $\Phi_{ij}\leq e^{\lambda_{\max}(M)}$, one has
$\norm{D\exp(M)[E]}_{\Frob}\leq e^{\lambda_{\max}(M)}\norm E_{\Frob}$. The Duhamel identity
$e^Y-e^X=\int_0^1\!\!\int_0^1e^{uM_s}(Y-X)e^{(1-u)M_s}\dd u\dd s$ with $M_s=X+s(Y-X)$, taken at
$X=\Delta tS$ and $Y=\Delta t\widehat S$ so that $M_s\preceq\Delta t\Id$ by convexity, together with
$\norm{e^{uM_s}}_\op\leq e^{u\Delta t}$, gives
$\norm{e^{\Delta t\widehat S}-e^{\Delta tS}}_{\Frob}\leq e^{\Delta t}\Delta t\norm\Xi_{\Frob}\leq e\Delta t\norm\Xi_{\Frob}$,
which is (i).

For (ii), Taylor's theorem with integral remainder gives
$e^{\Delta t\widehat S}-e^{\Delta tS}-D\exp(\Delta tS)[\Delta t\,\Xi]=\int_0^1(1-s)\Phi''(s)\dd s$ with
$\Phi(s):=\exp(\Delta tS+s\Delta t\,\Xi)$ and
$\Phi''(s)=D^2\exp(M_s)[\Delta t\,\Xi,\Delta t\,\Xi]$, $M_s\preceq\Delta t\Id$. The second Fr\'echet
derivative admits the simplex representation
$D^2\exp(M)[E,E]=2\int_{u_0+u_1+u_2=1,\,u_i\geq0}e^{u_0M}Ee^{u_1M}Ee^{u_2M}$, whose Frobenius norm is
bounded, using $\norm{e^{u_iM}}_\op\leq e^{u_i\lambda_{\max}(M)}$ and
$\norm{ABC}_{\Frob}\leq\norm A_\op\norm B_\op\norm C_{\Frob}$, by
$2\cdot\tfrac12\cdot e^{\lambda_{\max}(M)}\norm E_\op\norm E_{\Frob}$, the factor $\tfrac12$ being the
volume of the simplex. Hence
$\norm{\Phi''(s)}_{\Frob}\leq e^{\Delta t}\Delta t^2\norm\Xi_\op\norm\Xi_{\Frob}$, and
$\int_0^1(1-s)\dd s=\tfrac12$ with $e^{\Delta t}\leq e$ gives (ii).
\end{proof}

\begin{proof}[Proof of \Cref{lem:perturb-Riem}]
Let $a^+:=T_{\Riem}(a)$ and $\widehat a^+-a^+=(\delta_m,\delta_C)$. Expanding the whitened square,
\begin{equation}
\label{eq:perturb-expand}
    \E\bigl[\norm{\widehat a^+-a_\star}_\star^2\mid\Fil\bigr]
    =\norm{a^+-a_\star}_\star^2
    +2\inner{a^+-a_\star}{\E[\widehat a^+-a^+\mid\Fil]}_\star
    +\E\bigl[\norm{\widehat a^+-a^+}_\star^2\mid\Fil\bigr],
\end{equation}
and abbreviate $Q:=\E[\norm{\overline\varepsilon_b}^2+\tfrac12\norm{\overline\varepsilon_H}_{\Frob}^2\mid\Fil]$.
For the second moment, the mean components differ by $\delta_m=\Delta t\,C\,\overline\varepsilon_b$, so
$\E\norm{\delta_m}^2\leq\tfrac94\Delta t^2\E\norm{\overline\varepsilon_b}^2$. The covariance blocks are
$C^{1/2}e^{\Delta tS}C^{1/2}$ and $C^{1/2}e^{\Delta t\widehat S}C^{1/2}$ with
$S=\Id-C^{1/2}A(a)C^{1/2}\preceq\Id$, $\widehat S=\Id-C^{1/2}\overline HC^{1/2}\preceq\Id$, and
$\widehat S-S=-C^{1/2}\overline\varepsilon_HC^{1/2}=:\Xi_S$. By \Cref{lem:exp-perturb}(i) and
$\norm{\Xi_S}_{\Frob}\leq\tfrac32\norm{\overline\varepsilon_H}_{\Frob}$,
\[
    \norm{\delta_C}_{\Frob}\leq\norm C_\op\norm{e^{\Delta t\widehat S}-e^{\Delta tS}}_{\Frob}
    \leq\tfrac32\cdot e\Delta t\cdot\tfrac32\norm{\overline\varepsilon_H}_{\Frob}
    =\tfrac{9e}4\Delta t\norm{\overline\varepsilon_H}_{\Frob}.
\]
Therefore
$\E\norm{\widehat a^+-a^+}_\star^2\leq\tfrac94\Delta t^2\E\norm{\overline\varepsilon_b}^2+\tfrac{81e^2}{32}\Delta t^2\E\norm{\overline\varepsilon_H}_{\Frob}^2\leq\tfrac{81e^2}{16}\Delta t^2Q$,
using $\tfrac94\leq\tfrac{81e^2}{16}$.
The mean bias vanishes, $\E[\delta_m\mid\Fil]=0$. For the covariance,
$e^{\Delta t\widehat S}-e^{\Delta tS}=D\exp(\Delta tS)[\Delta t\,\Xi_S]+\mathcal R$ with the first term
centered, so $\E[\delta_C\mid\Fil]=C^{1/2}\E[\mathcal R\mid\Fil]C^{1/2}$ and, by
\Cref{lem:exp-perturb}(ii) with $\norm{\Xi_S}_\op\norm{\Xi_S}_{\Frob}\leq\tfrac94\norm{\overline\varepsilon_H}_{\Frob}^2$,
\[
    \norm{\E[\delta_C\mid\Fil]}_{\Frob}\leq\tfrac32\cdot\tfrac e2\Delta t^2\cdot\tfrac94\E\norm{\overline\varepsilon_H}_{\Frob}^2
    =\tfrac{27e}{16}\Delta t^2\E\norm{\overline\varepsilon_H}_{\Frob}^2 .
\]
Since the mean bias is zero, $\norm{\E[\widehat a^+-a^+\mid\Fil]}_\star=\tfrac1{\sqrt2}\norm{\E[\delta_C\mid\Fil]}_{\Frob}$,
so with $\norm{a^+-a_\star}_\star\leq1$, Cauchy--Schwarz, and
$\E\norm{\overline\varepsilon_H}_{\Frob}^2\leq2Q$, the cross term of \eqref{eq:perturb-expand} is at
most $\tfrac{27e}{4\sqrt2}\Delta t^2Q$.
Adding the two contributions, the coefficient of $\Delta t^2Q$ is
$\tfrac{81e^2}{16}+\tfrac{27e}{4\sqrt2}<64$, numerically about $50.4$. Bounding
$Q\leq B^{-1}(V_0+V_1\norm{a-a_\star}_\star^2)$ by \Cref{lem:loc-variance} gives the claim.
\end{proof}

\begin{proof}[Proof of \Cref{lem:perturb-KL}]
Let $a^+:=T_{\KL}(a)$, $P:=C^{-1}+\Delta t\,A(a)$ and
$\widehat P:=C^{-1}+\Delta t\,\overline H=P+\Delta t\,\overline\varepsilon_H$, so that
$C^+=(1+\Delta t)P^{-1}$, $\widehat C^+=(1+\Delta t)\widehat P^{-1}$ and
$\delta_C=(1+\Delta t)(\widehat P^{-1}-P^{-1})$. Since $A(a),\overline H\succeq0$ we have
$P\succeq C^{-1}$ and $\widehat P\succeq C^{-1}$, so
$\norm{P^{-1}}_\op,\norm{\widehat P^{-1}}_\op\leq\norm C_\op\leq\tfrac32$ almost surely. The expansion
\eqref{eq:perturb-expand} holds with $\norm\cdot_\star$ replaced by $\norm\cdot_{\star,\Delta t}$,
whose covariance weight is $\tfrac{1+\Delta t}2$.
For the second moment, as before $\E\norm{\delta_m}^2\leq\tfrac94\Delta t^2\E\norm{\overline\varepsilon_b}^2$.
The first-order resolvent identity
$\widehat P^{-1}-P^{-1}=-P^{-1}(\Delta t\,\overline\varepsilon_H)\widehat P^{-1}$ gives
$\norm{\widehat P^{-1}-P^{-1}}_{\Frob}\leq\tfrac94\Delta t\norm{\overline\varepsilon_H}_{\Frob}$, hence
$\norm{\delta_C}_{\Frob}\leq\tfrac92\Delta t\norm{\overline\varepsilon_H}_{\Frob}$. As
$\tfrac{1+\Delta t}2\leq1$, we get
$\E\norm{\widehat a^+-a^+}_{\star,\Delta t}^2\leq\tfrac94\Delta t^2\E\norm{\overline\varepsilon_b}^2+\tfrac{81}4\Delta t^2\E\norm{\overline\varepsilon_H}_{\Frob}^2\leq\tfrac{81}2\Delta t^2Q$.
For the bias, the second-order resolvent identity
\[
    \widehat P^{-1}-P^{-1}=-P^{-1}(\Delta t\,\overline\varepsilon_H)P^{-1}+P^{-1}(\Delta t\,\overline\varepsilon_H)P^{-1}(\Delta t\,\overline\varepsilon_H)\widehat P^{-1}
\]
has centered first term, so the conditional bias is the second term, and
$\norm{\E[\widehat P^{-1}-P^{-1}\mid\Fil]}_{\Frob}\leq(\tfrac32)^3\Delta t^2\E\norm{\overline\varepsilon_H}_{\Frob}^2=\tfrac{27}8\Delta t^2\E\norm{\overline\varepsilon_H}_{\Frob}^2$.
Thus $\norm{\E[\delta_C\mid\Fil]}_{\Frob}\leq\tfrac{27}4\Delta t^2\E\norm{\overline\varepsilon_H}_{\Frob}^2$
and, the mean bias being zero,
$\norm{\E[\widehat a^+-a^+\mid\Fil]}_{\star,\Delta t}\leq\tfrac{27}4\Delta t^2\E\norm{\overline\varepsilon_H}_{\Frob}^2$;
with $\norm{a^+-a_\star}_{\star,\Delta t}\leq1$ the cross term is at most $27\Delta t^2Q$.
Adding the two contributions, the coefficient of $\Delta t^2Q$ is $\tfrac{81}2+27=67.5<68$. Finally
$\norm{a-a_\star}_\star^2\leq\norm{a-a_\star}_{\star,\Delta t}^2$, so
$Q\leq B^{-1}(V_0+V_1\norm{a-a_\star}_\star^2)$ from \Cref{lem:loc-variance} gives the stated bound.
\end{proof}

\subsection{Killed-process contraction and finite-horizon containment}

\begin{proof}[Proof of \Cref{thm:stoch-loc-Riem}]
Write $r_n:=\norm{a_n-a_\star}_\star$. On $\{n<\tau_{\Riem}\}$ we have $a_n\in\Ureg_{\delta_{\Riem}}$,
so $\tfrac12\Id\preceq C_n\preceq\tfrac32\Id$ and
$r_n\leq\mathfrak b_{\star,\Riem}(\delta_{\Riem})\sqrt{\delta_{\Riem}}\leq1$ by
\Cref{def:stoch-admissible}. The pointwise one-step contraction of \Cref{lem:loc-contraction}, which
needs only $\Delta t\leq\tfrac2{4+\Gamma}$ and no energy-descent bound, gives
$\norm{T_{\Riem}(a_n)-a_\star}_\star\leq(1-\mu_{\Riem}\Delta t)r_n\leq r_n\leq1$, since
$\tfrac12\gamma_{\Riem}=\mu_{\Riem}$. Hence \Cref{lem:perturb-Riem} applies and
\[
    \E[r_{n+1}^2\mid\Fil_n]\leq(1-\mu_{\Riem}\Delta t)^2r_n^2+\tfrac{64\Delta t^2}B\bigl(V_0+V_1r_n^2\bigr).
\]
Because $\Delta t\leq\tfrac2{4+\Gamma}$ we have
$\mu_{\Riem}\Delta t=\tfrac{\Delta t}{2(4+\Gamma)}\leq\tfrac1{(4+\Gamma)^2}\leq\tfrac1{16}\leq\tfrac12$,
so $(1-\mu_{\Riem}\Delta t)^2\leq1-\tfrac32\mu_{\Riem}\Delta t$ from $x^2\leq\tfrac12x$ for
$x\leq\tfrac12$. The stepsize bound $\Delta t\leq\tfrac B{256A_{\Riem}V_1}$ gives
$\tfrac{64\Delta t^2}BV_1\leq\tfrac12\mu_{\Riem}\Delta t$, hence
\[
    \E[r_{n+1}^2\mid\Fil_n]\leq(1-\mu_{\Riem}\Delta t)r_n^2+\eta_{\Riem},
    \qquad \eta_{\Riem}:=\tfrac{64\Delta t^2}BV_0=\tfrac{192\Delta t^2N_\theta^2\LHs^2}B .
\]
Set $M_n:=\E[r_n^2\mathbf1_{\{\tau_{\Riem}>n\}}]$. Since
$\{\tau_{\Riem}>n+1\}\subseteq\{\tau_{\Riem}>n\}\in\Fil_n$,
\[
    M_{n+1}\leq\E\bigl[\mathbf1_{\{\tau_{\Riem}>n\}}\E[r_{n+1}^2\mid\Fil_n]\bigr]
    \leq(1-\mu_{\Riem}\Delta t)M_n+\eta_{\Riem},
\]
and iterating from $M_0=r_0^2$ with $\sum_{j\geq0}(1-\mu_{\Riem}\Delta t)^j=(\mu_{\Riem}\Delta t)^{-1}$
gives
$M_N\leq(1-\mu_{\Riem}\Delta t)^Nr_0^2+\eta_{\Riem}/(\mu_{\Riem}\Delta t)$, which is the stated bound
since $\eta_{\Riem}/(\mu_{\Riem}\Delta t)=384\Delta t(4+\Gamma)N_\theta^2\LHs^2/B$.
\end{proof}

\begin{proof}[Proof of \Cref{thm:stoch-loc-KL}]
Identical to the previous proof, replacing $\norm\cdot_\star$ by $\norm\cdot_{\star,\Delta t}$,
$\gamma_{\Riem}$ by $\gamma_{\KL}$ so that $\mu_{\Riem}$ becomes $\mu_{\KL}$,
\Cref{lem:perturb-Riem} by \Cref{lem:perturb-KL} with the coefficient $64$ replaced by $68$, and
$\eta_{\Riem}$ by $\eta_{\KL}=\tfrac{68\Delta t^2}BV_0=\tfrac{204\Delta t^2N_\theta^2\LHs^2}B$. The
step condition $\Delta t\leq\tfrac B{272A_{\KL}V_1}$ gives
$\tfrac{68\Delta t^2}BV_1\leq\tfrac12\mu_{\KL}\Delta t$, and $\mu_{\KL}\Delta t\leq\tfrac12$ holds as
before, so $\E[r_{n+1}^2\mid\Fil_n]\leq(1-\mu_{\KL}\Delta t)r_n^2+\eta_{\KL}$ with
$r_n=\norm{a_n-a_\star}_{\star,\Delta t}$; the non-exit recursion gives
$M_N\leq(1-\mu_{\KL}\Delta t)^Nr_0^2+\eta_{\KL}/(\mu_{\KL}\Delta t)$, which is the stated bound.
\end{proof}

\begin{proof}[Proof of \Cref{prop:exit-prob}]
Since $\{a:\norm{a-a_\star}_{\star,\bullet}\leq r_\bullet\}\subseteq\Ureg_{\delta_\bullet}$, on
$\{n<\tau_{r_\bullet}\}$ the one-step recursion
$\E[r_{n+1}^2\mid\Fil_n]\leq q_\bullet r_n^2+\eta_\bullet$ established in the two preceding proofs
holds, with $q_\bullet:=1-\mu_\bullet\Delta t$ and $r_n:=\norm{a_n-a_\star}_{\star,\bullet}$. Define the
process frozen at exit, $R_n:=r_{n\wedge\tau_{r_\bullet}}$. On $\{\tau_{r_\bullet}>n\}$ one has
$R_{n+1}=r_{n+1}$ and, since $q_\bullet\leq1$,
$\E[R_{n+1}^2\mid\Fil_n]\leq q_\bullet r_n^2+\eta_\bullet\leq R_n^2+\eta_\bullet$; on
$\{\tau_{r_\bullet}\leq n\}$ one has $R_{n+1}=R_n$. In either case
$\E[R_{n+1}^2\mid\Fil_n]\leq R_n^2+\eta_\bullet$, so iterating gives $\E R_N^2\leq r_0^2+N\eta_\bullet$.
On $\{\tau_{r_\bullet}\leq N\}$ the frozen value satisfies $R_N=r_{\tau_{r_\bullet}}>r_\bullet$, so
Markov's inequality gives
$r_\bullet^2\Prob(\tau_{r_\bullet}\leq N)\leq\E R_N^2\leq r_0^2+N\eta_\bullet$, which is the claim.
Splitting the failure budget equally between the two terms gives the displayed sufficient conditions.
\end{proof}

Freezing the process at the exit time removes the factor $(\mu_\bullet\Delta t)^{-1}$ that a
per-exit-time union bound would incur; that denominator remains genuine in the stationary killed-moment
floor $\eta_\bullet/(\mu_\bullet\Delta t)$ of \Cref{thm:stoch-loc-Riem,thm:stoch-loc-KL}, which sums
repeated noise within the surviving process, but not in the exit probability. At the certified stepsize
$\Delta t=\Theta((4+\Gamma)^{-1})$, where also $\mu_\bullet=\Theta((4+\Gamma)^{-1})$, the resulting
batch requirement is smaller by a factor of order $(4+\Gamma)^2$ than the condition
$B\gtrsim N\Delta t\,N_\theta^2\LHs^2/(\epsilon'\mu_\bullet r_\bullet^2)$ that the union-bound argument
would demand.

\begin{proof}[Proof of \Cref{cor:stoch-three-stage}]
By Markov's inequality applied to $\Delta_{N_1}\geq0$ and $\Ureg_{\delta_{\mathrm{in}}}=\{\Delta\leq\delta_{\mathrm{in}}\}$,
\[
    \Prob\bigl(a_{N_1}\notin\Ureg_{\delta_{\mathrm{in}}}\bigr)=\Prob(\Delta_{N_1}>\delta_{\mathrm{in}})
    \leq\frac{\E\Delta_{N_1}}{\delta_{\mathrm{in}}}\leq\epsilon_{\mathrm{ent}} .
\]
On $E_{\mathrm{in}}$ the norm conversion \eqref{eq:band-conversion} gives
$\norm{a_{N_1}-a_\star}_{\star,\bullet}\leq\mathfrak b_{\star,\bullet}(\delta_{\mathrm{in}})\sqrt{\delta_{\mathrm{in}}}=:r_0$,
and the deep-entry choice makes $r_0^2\leq\tfrac{\epsilon'}2r_\bullet^2\leq r_\bullet^2$, so $a_{N_1}$
lies in the outer ball. Since $E_{\mathrm{in}}\in\Fil_{N_1}$, applying \Cref{prop:exit-prob} to the
process $(a_{N_1+n})_{n\geq0}$ on $E_{\mathrm{in}}$, the deep-entry condition controls the
initial-condition term by $\epsilon'/2$ and the batch condition controls the $N_2\eta_\bullet$ term by
$\epsilon'/2$, giving $\Prob(\{\tau^{(N_1)}\leq N_2\}\cap E_{\mathrm{in}})\leq\epsilon'$. Hence
$\Prob(E_{\mathrm{in}}\cap\{\tau^{(N_1)}>N_2\})\geq(1-\epsilon_{\mathrm{ent}})-\epsilon'$.

For the second moment write $\rho_n:=\norm{a_{N_1+n}-a_\star}_{\star,\bullet}$. On
$\{n<\tau^{(N_1)}\}\cap E_{\mathrm{in}}$ the outer ball contains $a_{N_1+n}$, so the one-step recursion
$\E[\rho_{n+1}^2\mid\Fil_{N_1+n}]\leq q_\bullet\rho_n^2+\eta_\bullet$ holds there; setting
$M_n:=\E[\rho_n^2\mathbf1_{E_{\mathrm{in}}}\mathbf1_{\{\tau^{(N_1)}>n\}}]$ and using
$\{\tau^{(N_1)}>n\}\cap E_{\mathrm{in}}\in\Fil_{N_1+n}$ gives $M_{n+1}\leq q_\bullet M_n+\eta_\bullet$,
so $M_{N_2}\leq q_\bullet^{N_2}r_0^2+\eta_\bullet/(\mu_\bullet\Delta t)$, and
$\eta_\bullet/(\mu_\bullet\Delta t)$ equals the displayed $F_{\Riem}$ resp.\ $F_{\KL}$.
\end{proof}

\subsection{The decreasing schedule}

\begin{proof}[Proof of \Cref{thm:decreasing}]
Since $n+n_0\geq n_0\geq64\kappa^2$, we have $\Delta t_n=8\kappa/(n+n_0)\leq1/(8\kappa)$, and the
covariance bands persist along the varying schedule. For the KL scheme this follows at each step from
the induction of \Cref{lem:stoch-band}(2), which allows every positive stepsize. For the Riemannian
scheme, suppose $\tfrac1{2\beta}\Id\preceq C_n\preceq\tfrac1\alpha\Id$; since $\Delta t_n\kappa\leq1$
and $\Delta t_n\leq1$, the upper-band argument gives
$C_{n+1}^{-1}\succeq(1-\Delta t_n)C_n^{-1}+\Delta t_n\overline H_n\succeq\alpha\Id$, and the lower
envelope gives
\[
    C_{n+1}^{-1}\preceq e^{-\Delta t_n}\bigl[C_n^{-1}+(e-1)\Delta t_n\overline H_n\bigr]
    \preceq\beta e^{-\Delta t_n}\bigl[2+(e-1)\Delta t_n\bigr]\Id\preceq2\beta\Id,
\]
because $e^{-h}[2+(e-1)h]\leq2$ for $0\leq h\leq1$. The rescued initializer lies in both bands, so
induction proves their persistence and the constant-stepsize one-step estimates apply at each $n$.

For the Riemannian scheme, \Cref{lem:onestep-stoch-Riem} with $\Delta t_n\leq1/(8\kappa)$ has bracket
$1-\kappa\Delta t_n(1+2/B)\geq1-\tfrac38\geq\tfrac12$, so the drift is at least
$\tfrac{\Delta t_n}2\Gradnorm^2$, and the Riemannian gradient domination
$\Gradnorm^2\geq\tfrac1{2\kappa}\Delta$ of \eqref{eq:PL-bands} makes it at least
$\tfrac{\Delta t_n}{4\kappa}\Delta_n$, with floor coefficient $\tfrac32\leq c_{\mathsf S}$. For the KL
scheme, \Cref{lem:onestep-stoch-KL} gives drift
$\tfrac{\Delta t_n}2\Gradnorm^2\geq\tfrac{\Delta t_n}{2\kappa}\Delta_n\geq\tfrac{\Delta t_n}{4\kappa}\Delta_n$
and floor coefficient $\tfrac74=c_{\mathsf S}$. In both cases, with $\Psi\leq\Psi_H$ from
\Cref{lem:Psi-hesslip},
\begin{equation}
\label{eq:decreasing-recursion}
    \E\Delta_{n+1}\leq\Bigl(1-\frac{\Delta t_n}{4\kappa}\Bigr)\E\Delta_n+\frac{c_{\mathsf S}\kappa\Delta t_n^2}B\Psi_H .
\end{equation}
Write $t:=n+n_0$, so that $\tfrac{\Delta t_n}{4\kappa}=\tfrac2t$ and
$\tfrac{c_{\mathsf S}\kappa\Delta t_n^2}B\Psi_H=\tfrac S{t^2}$ with
$S:=64c_{\mathsf S}\kappa^3\Psi_H/B$, and \eqref{eq:decreasing-recursion} becomes
$\E\Delta_{n+1}\leq(1-\tfrac2t)\E\Delta_n+\tfrac S{t^2}$. We show $\E\Delta_n\leq K/t$ by induction
with $K=\max\{S,n_0\Delta(a_0)\}$. For $n=0$, $\E\Delta_0=\Delta(a_0)\leq K/n_0$. Assuming
$\E\Delta_n\leq K/t$,
\[
    \E\Delta_{n+1}\leq\Bigl(1-\frac2t\Bigr)\frac Kt+\frac S{t^2}
    =\frac Kt-\frac{2K-S}{t^2}
    \leq\frac Kt-\frac K{t^2}
    =\frac{K(t-1)}{t^2}
    \leq\frac K{t+1},
\]
using $2K-S\geq K$ and $(t-1)(t+1)\leq t^2$. Since $t+1=(n+1)+n_0$, the induction closes.
\end{proof}

\begin{remark}[Conservative constants]
\label{rem:stoch-constants}
The explicit constants entering the local floors, namely $V_0=3N_\theta^2\LHs^2$ and
$V_1=24+6N_\theta\LHs^2$ of \Cref{lem:loc-variance}, the aggregate coefficients $64$ and $68$ of
\Cref{lem:perturb-Riem,lem:perturb-KL}, the exponential remainder constant $e/2$ of
\Cref{lem:exp-perturb}, the stepsize denominators $256$ and $272$, and the floor coefficients $192$,
$204$, $384$, $408$, are conservative and not optimized. Only their orders in
$(\Delta t,N_\theta,\LHs,B,A_\bullet)$ enter the conclusions of \Cref{sec:stochastic}. The same applies
to the global constants: the retraction-smoothness constant $L=2\beta\Lambda$ may be replaced by
$2(e-2)\beta\Lambda\approx1.437\beta\Lambda$, which rescales the stepsizes and floors of
\Cref{thm:stoch-glob-Riem,thm:stoch-glob-KL} by the factor $e-2$ without affecting any order, and the
dimension coefficient $2\beta\Lambda$ of \Cref{lem:Psi-selfbound} may be reduced to at most
$\tfrac54\beta\Lambda$ by retaining the centering term discarded in its first step.
\end{remark}

% SOURCES:
% improved-global-local-stoch.tex l.2258-2323 (proof of lem:rescue-energy) -> proof of lem:rescue-energy
% improved-global-local-stoch.tex l.2546-2558 (proof of lem:Riem-onestep-stoch) -> proof of lem:onestep-stoch-Riem
% improved-global-local-stoch.tex l.2690-2767 (proof of lem:KL-onestep-stoch) -> proof of lem:onestep-stoch-KL
% improved-global-local-stoch.tex l.2575-2582 (proof of thm:Riem-Hlip) -> proof of thm:stoch-glob-Riem(1)
% improved-global-local-stoch.tex l.2601-2617 (proof of thm:Riem-selfbound) -> proof of thm:stoch-glob-Riem(2)
% improved-global-local-stoch.tex l.2783-2791 (proof of thm:KL-Hlip) -> proof of thm:stoch-glob-KL(1)
% improved-global-local-stoch.tex l.2827-2842 (proof of thm:KL-selfbound) -> proof of thm:stoch-glob-KL(2)
% improved-global-local-stoch.tex l.2994-3058 (proof of lem:local-variance) -> proof of lem:loc-variance
% improved-global-local-stoch.tex l.3066-3104 (lem:exp-perturb and proof) -> lem:exp-perturb
%   (source symbol \Delta for the increment renamed to \Xi to avoid collision with the energy gap)
% improved-global-local-stoch.tex l.3118-3171 (proof of lem:local-perturb-Riem) -> proof of lem:perturb-Riem
% improved-global-local-stoch.tex l.3184-3221 (proof of lem:local-perturb-KL) -> proof of lem:perturb-KL
% improved-global-local-stoch.tex l.3250-3281 (proof of thm:stoch-local-Riem) -> proof of thm:stoch-loc-Riem
% improved-global-local-stoch.tex l.3303-3313 (proof of thm:stoch-local-KL) -> proof of thm:stoch-loc-KL
% improved-global-local-stoch.tex l.3337-3366, l.3368-3386 (prop:finite-horizon and proof) -> proof of prop:exit-prob
%   and the following paragraph on the frozen-process gain
% improved-global-local-stoch.tex l.3579-3608 (proof of cor:stoch-three-stage) -> proof of cor:stoch-three-stage
% improved-global-local-stoch.tex l.3675-3721 (proof of thm:decreasing) -> proof of thm:decreasing
% improved-global-local-stoch.tex l.3628-3631, l.3916-3918 (conservative-constant notes),
%   l.1003-1005 (rem:retr-sharp), l.2455-2466 (rem:selfbound-sharp) -> rem:stoch-constants
